\scp{Possibilities for combining some Wallacean subgroups}{13-04}
In chapters \ref{ch:BL}--\ref{ch:Ban} we explored the unique historical phonologies
of individual Wallacean subgroups,
along with internal divisions within them.
In \srf{02-02} and \srf{02-03}, we described various kinds of
evidence showing extensive movement
and trade that has spanned two or more millennia,
as well as contact with typologically different Papuan languages
spanning almost four millennia within the Wallacean region,
which has been intense enough to not only result in borrowed
vocabulary, but also resulted in uneven changes in the grammatical
and typological structures of many Austronesian languages across the region.
In \srf{13-01}, \srf{13-02}, and \srf{13-03}, we took a more
bird’s-eye view across the subgroups,
which revealed some similarities in the historical phonologies,
phonotactics and morphosyntax between a few Wallacean subgroups,
but none that rose to the level of subgrouping arguments that
allow us to join them together with any confidence.
We revisit a few additional combinations in this brief section.

Adjacent subgroups may share some features in the phonology
and lexicon that are not shared by non-adjacent subgroups,
hence their description as a “chain” (see \crf{ch:CM},
for example, although some of those leap over adjacent groups
as is also common in chaining).
First we revisit some general issues,
and review drawbacks for basing subgrouping on those.

\begin{itemize}
	\item	Merger of PMP *ñ/*n > \it{n}: In \srf{13-01-01} it was noted
				that: (a) evidence for any given language for reflexes of PMP
				*ñ can be elusive; (b) individual languages in several subgroups
				show sporadic raising of vowels adjacent to PMP *ñ,
				leaving traces such that the merger of *ñ/*n > \it{n} for those
				subgroups may not be correct; (c) the merger of *ñ/*n
				> \it{n} is common cross-linguistically; (d) the merger of *ñ/*n
				> \it{n} is also found outside the target region;
				and (e) *ñ is retained unchanged in \tsc{\ili{Bima}-Lembata} (specifically,
				\tsc{\ili{Havu}-Dhao}) 
	\item	Weakening of PMP *p > \it{ɸ,
				f, h,} {\0}: In \srf{13-01-02} it was noted that: (a) within
				those Wallacean subgroups that show such weakening,
				there are some languages that retain *p = \it{p} in (at l\ili{east}
				some words in) the initial,
				medial or final position.
				So strictly speaking some could argue that those subgroups should
				be described as retaining PMP *p = \it{p},
				at l\ili{east} for those positions
				and words in which it is retained in some languages; (b) reconciling
				the discrepancies between subgroups that retain some initial
				*p- with those that retain some final *-p would require complex
				caveats, unless one made the broader generalisation that they all retain
				*p = \it{p}, which then disqualifies it as a subgrouping argument; (c) a weakening
				of *p to a voiceless fricative is fairly common cross-linguistically;
				and (d) a weakening of PMP *p is also found outside the region,
				such as in SHWNG.
\end{itemize}

Barring more convincing evidence we are hesitant to join subgroups
on such a weak basis which could easily be due to convergence.
The possible combinations we explored include:

\begin{itemize}
	\item	\tsc{Sula-Buru} + \tsc{\ili{Ambon}-Seram} = \ili{Central} Maluku: *a/*-ay
				> \it{a}; *n/*ñ > \it{n}.
				Both mergers are extremely common cross-linguistically,
				within the region, and are not exclusive.
				Uncommon sound change(s) that are exclusive would be more persuasive.
				See \crf{ch:CM} for the bundles of features that distinguish
				these two subgroups, including in the morphosyntax.
				An earlier draft of this book treated these two as a single subgroup,
				but the extensive and systematic differences are far more significant
				than the very weak “innovations” they share which are not exclusive
				and could easily be from parallel development.
				It is possible that they were together briefly
				and split at a very early stage,
				but their many differences suggest otherwise.
	\item	\tsc{\ili{Ambon}-Seram} + \tsc{Seram-Tanimbar-Bomberai} = Maluku-Bomberai:
				*a/\mbox{*-ay} > \it{a}; *d/*z > **d: These two mergers together may
				be a possibility, although the first is weak and not exclusive.
				While both subgroups merge PMP *d/*z,
				the reflexes are different.
				In \tsc{\ili{Ambon}-Seram} *d/*z > **r,
				while in STB *d/*z > **d,
				at l\ili{east} word initially (see \srf{07-03-03}).
				We find the behaviour of the rest of the historical apicals however,
				to be difficult to reconcile as coming from a common parent language.
				So for the moment,
				we see the merger of PMP *d/*z as parallel development.
	\item	\ili{Banda} + \tsc{\ili{Ambon}-Seram}? \ili{Banda} + \tsc{Seram-Tanimbar-Bomberai}?:
				\ili{Banda} shares some features in its historical phonology with both
				subgroups, and is also divergent from both (see \crf{ch:Ban}).
				The shared mergers of PMP nasals *ŋ/*n/*ñ > \it{n}
				and *d/*z > \it{r} could link it with \tsc{\ili{Ambon}-Seram},
				but the different mergers involving *R,
				and the different behaviour of *a make such a combination difficult.
				Lexically, \ili{Banda} is equally (un)related to languages in both
				subgroups, as well as with some \tsc{Timor-Babar} languages \citep{Hughes1987,LoskiLoski1989}.
				It is also lexically
				and morphologically divergent from the languages with which it
				was proposed to belong (see \citealp[141ff]{Collins1986}).
				Having been displaced from its traditional lands for four centuries
				surely complicates the picture.
				There is no agreement among investigators as to how \ili{Banda} should be classified.
	\item	\tsc{Sula-Buru} + \tsc{Seram-Tanimbar-Bomberai}: *a/*-ay/*-aw
				> \it{a}: \citet[264--266]{Blust1993} saw “glide truncation” as one
				of the most “distinctive features” of CMP,
				so this pattern is more widespread than these two subgroups (\srf{13-01-03}).
				The very different behaviour of the PMP labials
				and apicals make the historical phonologies of these two subgroups difficult to reconcile.
\end{itemize}

\ili{As} noted in \crf{ch:LayAdd},
in trying to see where the linguistic picture
and the archaeological picture link up,
\citet[511]{Spriggs2011} comments, “We can use the spread of
the ISEA Neolithic as a proxy for AN language spread,
as justified at length by \citet{Pawley2004}
and \citet{Ross2008}, among others.
\it{In doing this, it is very hard to see anything between PMP
and EMP at all from the archaeology}” [emphasis ours].
\ili{As} yet there appears to be no evidence in the archaeological
record that mirrors a pause or a node at PCEMP.

We have, as yet,
been unable to find anything in the historical phonology or morphosyntax
from within the Wallacean region that would support either of
the purported CMP or CEMP subgroups.
That leaves the lexicon as perhaps the strongest basis for the CEMP hypothesis.
Yet as we have been showing in this zone of long-term contact,
mobility and trade, basing a subgroup on lexical data is highly problematic.
We explore this further in the next two chapters.